\section{The large-neighbourhood branch}

\import{function.tex}
\import{combinatorics/graph-cardinality.tex}
\import{set/equinumerosity.tex}
\import{set/iteration.tex}
\import{set/choice.tex}
\import{cardinal.tex}
\import{cardinal/countability.tex}

\subsection{Large neighbourhoods and good vertices}

\begin{definition}\label{large_neighbourhood}
    $\largeNeighbourhood{E}{X}{\kappa}$ iff for every $Y$ such that
        $Y\subseteq X$ and $Y\approx\kappa$ there exists $x\in Y$ such that
        $\graphNeighbourhood{E}{x}{Y}\approx\kappa$.
\end{definition}

\begin{definition}\label{edm_good_vertices}
    $\edmGoodVertices{E}{Y}{\kappa}
        =\{x\in Y\mid
            \graphNeighbourhood{E}{x}{Y}\approx\kappa\}$.
\end{definition}

\begin{proposition}\label{edm_good_vertices_iff}
    $x\in\edmGoodVertices{E}{Y}{\kappa}$ iff
        $x\in Y$ and
        $\graphNeighbourhood{E}{x}{Y}\approx\kappa$.
\end{proposition}
\begin{proof}
    Follows by \obv{edm_good_vertices}.
\end{proof}

\begin{proposition}\label{edm_good_vertices_inhabited}
    Suppose $\largeNeighbourhood{E}{X}{\kappa}$.
    Suppose $Y\subseteq X$ and $Y\approx\kappa$.
    Then $\edmGoodVertices{E}{Y}{\kappa}$ is inhabited.
\end{proposition}
\begin{proof}
    Take $x\in Y$ such that
        $\graphNeighbourhood{E}{x}{Y}\approx\kappa$ by
        \obv{large_neighbourhood}.
    Hence $x\in\edmGoodVertices{E}{Y}{\kappa}$ by
        \obv{edm_good_vertices_iff}.
\end{proof}

\begin{proposition}\label{edm_good_vertex_choice}
    Suppose $\largeNeighbourhood{E}{X}{\kappa}$.
    Suppose $Y\subseteq X$ and $Y\approx\kappa$.
    Then $\choice{\edmGoodVertices{E}{Y}{\kappa}}
        \in\edmGoodVertices{E}{Y}{\kappa}$.
\end{proposition}
\begin{proof}
    $\edmGoodVertices{E}{Y}{\kappa}$ is inhabited by
        \obv{edm_good_vertices_inhabited}.
    Follows by \obv{choice_in}.
\end{proof}

\subsection{The reservoir step}

\begin{definition}\label{edm_reservoir_step}
    $\edmReservoirStep{E}{X}{\kappa}
        =\{(Y,
            \graphNeighbourhood{E}
                {\choice{\edmGoodVertices{E}{Y}{\kappa}}}{Y})
            \mid Y\in\pow{X}\}$.
\end{definition}

\begin{theorem}\label{edm_reservoir_step_type}
    $\edmReservoirStep{E}{X}{\kappa}
        \in\funs{\pow{X}}{\pow{X}}$.
\end{theorem}
\begin{proof}
    Let $T=\edmReservoirStep{E}{X}{\kappa}$.
    $T$ is a relation by \auto{edm_reservoir_step,relation}.
    $T$ is right-unique by
        \auto{edm_reservoir_step,rightunique,pair_eq_iff}.
    $\dom{T}=\pow{X}$ by
        \auto{setext,dom_iff,edm_reservoir_step,pair_eq_iff}.
    Show that for every $Y\in\pow{X}$ we have
        $\graphNeighbourhood{E}
            {\choice{\edmGoodVertices{E}{Y}{\kappa}}}{Y}
            \in\pow{X}$.
    \begin{subproof}
        Fix $Y\in\pow{X}$.
        Hence $Y\subseteq X$ by \obv{powerset_elim}.
        $\graphNeighbourhood{E}
            {\choice{\edmGoodVertices{E}{Y}{\kappa}}}{Y}
            \subseteq Y$ by \obv{graph_neighbourhood_subseteq}.
        Hence $\graphNeighbourhood{E}
            {\choice{\edmGoodVertices{E}{Y}{\kappa}}}{Y}
            \subseteq X$ by \obv{subseteq_transitive}.
        Follows by \obv{powerset_intro}.
    \end{subproof}
    For every $Y\in\dom{T}$ we have
        $T(Y)=\graphNeighbourhood{E}
            {\choice{\edmGoodVertices{E}{Y}{\kappa}}}{Y}$ by
        \auto{edm_reservoir_step,function_apply_intro}.
    Hence $T$ is a function to $\pow{X}$ by \auto.
    Follows by \auto{funs_intro}.
\end{proof}

\begin{proposition}\label{edm_reservoir_step_apply}
    Let $Y\in\pow{X}$.
    Then $(\edmReservoirStep{E}{X}{\kappa})(Y)
        =\graphNeighbourhood{E}
            {\choice{\edmGoodVertices{E}{Y}{\kappa}}}{Y}$.
\end{proposition}
\begin{proof}
    $\edmReservoirStep{E}{X}{\kappa}$ is a function by
        \auto{edm_reservoir_step_type,funs_is_function}.
    $(Y,\graphNeighbourhood{E}
            {\choice{\edmGoodVertices{E}{Y}{\kappa}}}{Y})
        \in\edmReservoirStep{E}{X}{\kappa}$ by
        \obv{edm_reservoir_step}.
    Follows by \obv{function_apply_intro}.
\end{proof}

\subsection{The reservoir iteration}

\begin{proposition}\label{edm_reservoir_iteration_type}
    $\iterationSet{\edmReservoirStep{E}{X}{\kappa}}{X}
        \in\funs{\naturals}{\pow{X}}$.
\end{proposition}
\begin{proof}
    $\edmReservoirStep{E}{X}{\kappa}
        \in\funs{\pow{X}}{\pow{X}}$ by
        \obv{edm_reservoir_step_type}.
    $X\in\pow{X}$ by \obv{powerset_intro,subseteq_refl}.
    Follows by \obv{iteration_set_type}.
\end{proof}

\begin{proposition}\label{edm_reservoir_iteration_zero}
    $(\iterationSet{\edmReservoirStep{E}{X}{\kappa}}{X})(\emptyset)=X$.
\end{proposition}
\begin{proof}
    $\edmReservoirStep{E}{X}{\kappa}
        \in\funs{\pow{X}}{\pow{X}}$ by
        \obv{edm_reservoir_step_type}.
    $X\in\pow{X}$ by \obv{powerset_intro,subseteq_refl}.
    Follows by \obv{iteration_set_zero}.
\end{proof}

\begin{proposition}\label{edm_reservoir_iteration_successor}
    Let $n$ be a natural number.
    Then
    \[
        (\iterationSet{\edmReservoirStep{E}{X}{\kappa}}{X})(\suc{n})
        =\graphNeighbourhood{E}
            {\choice{\edmGoodVertices{E}
                {(\iterationSet{\edmReservoirStep{E}{X}{\kappa}}{X})(n)}
                {\kappa}}}
            {(\iterationSet{\edmReservoirStep{E}{X}{\kappa}}{X})(n)}
    \].
\end{proposition}
\begin{proof}
    Let $T=\edmReservoirStep{E}{X}{\kappa}$.
    Let $R=\iterationSet{\edmReservoirStep{E}{X}{\kappa}}{X}$.
    $T\in\funs{\pow{X}}{\pow{X}}$ by
        \auto{edm_reservoir_step_type}.
    $X\in\pow{X}$ by \obv{powerset_intro,subseteq_refl}.
    $R\in\funs{\naturals}{\pow{X}}$ by
        \auto{edm_reservoir_iteration_type}.
    Hence $R(n)\in\pow{X}$ by \obv{funs_type_apply}.
    $R(\suc{n})=T(R(n))$ by
        \auto{iteration_set_successor}.
    $T(R(n))=\graphNeighbourhood{E}
        {\choice{\edmGoodVertices{E}{R(n)}{\kappa}}}{R(n)}$ by
        \auto{edm_reservoir_step_apply}.
\end{proof}

\subsection{A countably infinite clique}

\begin{theorem}\label{large_neighbourhood_gives_clique}
    Let $\kappa$ be an infinite cardinal.
    Suppose $X\approx\kappa$.
    Suppose $E$ is a simple graph on $X$.
    Suppose $\largeNeighbourhood{E}{X}{\kappa}$.
    Then $E$ has a countably infinite clique in $X$.
\end{theorem}
\begin{proof}
    Let $T=\edmReservoirStep{E}{X}{\kappa}$.
    Let $R=\iterationSet{\edmReservoirStep{E}{X}{\kappa}}{X}$.
    $T\in\funs{\pow{X}}{\pow{X}}$ by
        \auto{edm_reservoir_step_type}.
    $R\in\funs{\naturals}{\pow{X}}$ by
        \auto{edm_reservoir_iteration_type}.
    Hence $R$ is a function and $\dom{R}=\naturals$ by
        \obv{funs_is_function,funs}.

    Let $v=\{(n,\choice{\edmGoodVertices{E}{R(n)}{\kappa}})
        \mid n\in\naturals\}$.
    $v$ is a relation by \auto{relation}.
    $v$ is right-unique by \auto{rightunique,pair_eq_iff}.
    $\dom{v}=\naturals$ by \auto{setext,dom_iff,pair_eq_iff}.
    Hence $v$ is a function on $\naturals$.
    For every $n\in\naturals$ we have
        $v(n)=\choice{\edmGoodVertices{E}{R(n)}{\kappa}}$ by
        \obv{function_apply_intro}.

    Let $C=\{n\in\naturals\mid
        R(n)\subseteq X\land
        R(n)\approx\kappa\land
        v(n)\in\edmGoodVertices{E}{R(n)}{\kappa}\land
        R(\suc{n})=\graphNeighbourhood{E}{v(n)}{R(n)}\}$.
    Show that $\emptyset\in C$.
    \begin{subproof}
        $\emptyset\in\naturals$ by \obv{omega_empty}.
        $R(\emptyset)=X$ by
            \auto{edm_reservoir_iteration_zero}.
        Hence $R(\emptyset)\subseteq X$ and
            $R(\emptyset)\approx\kappa$ by
            \obv{subseteq_refl}.
        Hence $\edmGoodVertices{E}{R(\emptyset)}{\kappa}$ is inhabited by
            \obv{edm_good_vertices_inhabited}.
        Hence $\choice{\edmGoodVertices{E}{R(\emptyset)}{\kappa}}
            \in\edmGoodVertices{E}{R(\emptyset)}{\kappa}$ by
            \obv{choice_in}.
        $v(\emptyset)=
            \choice{\edmGoodVertices{E}{R(\emptyset)}{\kappa}}$.
        Hence $v(\emptyset)\in
            \edmGoodVertices{E}{R(\emptyset)}{\kappa}$.
        \begin{align*}
            R(\suc{\emptyset})
                &= \graphNeighbourhood{E}
                    {\choice{\edmGoodVertices{E}{R(\emptyset)}{\kappa}}}
                    {R(\emptyset)}
                    \explanation{by
                        \auto{edm_reservoir_iteration_successor}}
                \\
                &= \graphNeighbourhood{E}{v(\emptyset)}{R(\emptyset)}
        \end{align*}
    \end{subproof}
    Show that for every $n\in C$ we have $\suc{n}\in C$.
    \begin{subproof}
        Fix $n\in C$.
        Hence $n\in\naturals$ and
            $R(n)\subseteq X$ and
            $R(n)\approx\kappa$ and
            $v(n)\in\edmGoodVertices{E}{R(n)}{\kappa}$ and
            $R(\suc{n})=\graphNeighbourhood{E}{v(n)}{R(n)}$.
        $\suc{n}\in\naturals$ by \obv{omega_successor}.
        Hence $R(\suc{n})\in\pow{X}$ by \obv{funs_type_apply}.
        Hence $R(\suc{n})\subseteq X$ by \obv{powerset_elim}.
        $\graphNeighbourhood{E}{v(n)}{R(n)}\approx\kappa$ by
            \auto{edm_good_vertices_iff}.
        Hence $R(\suc{n})\approx\kappa$.
        Hence $\edmGoodVertices{E}{R(\suc{n})}{\kappa}$ is inhabited by
            \obv{edm_good_vertices_inhabited}.
        Hence $\choice{\edmGoodVertices{E}{R(\suc{n})}{\kappa}}
            \in\edmGoodVertices{E}{R(\suc{n})}{\kappa}$ by
            \obv{choice_in}.
        $v(\suc{n})=
            \choice{\edmGoodVertices{E}{R(\suc{n})}{\kappa}}$.
        Hence $v(\suc{n})\in
            \edmGoodVertices{E}{R(\suc{n})}{\kappa}$.
        \begin{align*}
            R(\suc{\suc{n}})
                &= \graphNeighbourhood{E}
                    {\choice{\edmGoodVertices{E}{R(\suc{n})}{\kappa}}}
                    {R(\suc{n})}
                    \explanation{by
                        \auto{edm_reservoir_iteration_successor}}
                \\
                &= \graphNeighbourhood{E}{v(\suc{n})}{R(\suc{n})}
        \end{align*}
    \end{subproof}
    Hence $C$ is an inductive set by \obv{inductive_set}.
    Hence $\naturals\subseteq C$ by \obv{omega_minimal,subseteq}.
    Show that for every $n\in\naturals$ we have $v(n)\in X$.
    \begin{subproof}
        Fix $n\in\naturals$.
        Hence $n\in C$ by \obv{subseteq}.
        Hence $R(n)\subseteq X$ and
            $v(n)\in\edmGoodVertices{E}{R(n)}{\kappa}$.
        Hence $v(n)\in R(n)$ by \obv{edm_good_vertices_iff}.
        Follows by \obv{subseteq}.
    \end{subproof}
    Hence $v$ is a function to $X$.
    Hence $v\in\funs{\naturals}{X}$ by \obv{funs_intro}.

    Let $D=\{n\in\naturals\mid
        \text{for every $m\in n$ we have
            $R(n)\subseteq\graphNeighbourhood{E}{v(m)}{X}$}\}$.
    Show that $\emptyset\in D$.
    \begin{subproof}
        $\emptyset\in\naturals$ by \obv{omega_empty}.
        For every $m\in\emptyset$ we have
            $R(\emptyset)\subseteq\graphNeighbourhood{E}{v(m)}{X}$ by
            \auto{emptyset}.
    \end{subproof}
    Show that for every $n\in D$ we have $\suc{n}\in D$.
    \begin{subproof}
        Fix $n\in D$.
        Hence $n\in\naturals$ and for every $m\in n$ we have
            $R(n)\subseteq\graphNeighbourhood{E}{v(m)}{X}$.
        Hence $n\in C$ by \obv{subseteq}.
        Hence $R(n)\subseteq X$ and
            $R(\suc{n})=\graphNeighbourhood{E}{v(n)}{R(n)}$.
        $\suc{n}\in\naturals$ by \obv{omega_successor}.
        Show that for every $m\in\suc{n}$ we have
            $R(\suc{n})\subseteq\graphNeighbourhood{E}{v(m)}{X}$.
        \begin{subproof}
            Fix $m\in\suc{n}$.
            \begin{byCase}
                \caseOf{$m=n$.}
                    Hence $\graphNeighbourhood{E}{v(n)}{R(n)}
                        \subseteq\graphNeighbourhood{E}{v(n)}{X}$ by
                        \obv{graph_neighbourhood_monotone}.
                    $R(\suc{n})=
                        \graphNeighbourhood{E}{v(n)}{R(n)}$.
                \caseOf{$m\neq n$.}
                    Hence $m\in n$ by \auto{suc_iff}.
                    Hence $R(n)\subseteq
                        \graphNeighbourhood{E}{v(m)}{X}$.
                    $R(\suc{n})=
                        \graphNeighbourhood{E}{v(n)}{R(n)}$.
                    $\graphNeighbourhood{E}{v(n)}{R(n)}\subseteq R(n)$ by
                        \obv{graph_neighbourhood_subseteq}.
                    Hence $R(\suc{n})\subseteq R(n)$.
                    Show that every element of $R(\suc{n})$ is an element of
                        $\graphNeighbourhood{E}{v(m)}{X}$.
                    \begin{subproof}
                        Fix $z\in R(\suc{n})$.
                        Hence $z\in R(n)$ by \obv{subseteq}.
                        Follows by \obv{subseteq}.
                    \end{subproof}
                    Follows by \obv{subseteq}.
            \end{byCase}
        \end{subproof}
    \end{subproof}
    Hence $D$ is an inductive set by \obv{inductive_set}.
    Hence $\naturals\subseteq D$ by \obv{omega_minimal,subseteq}.

    Show that $v$ is injective.
    \begin{subproof}
        Show that for all $m,n$ such that $m,n\in\dom{v}$ and
            $v(m)=v(n)$ we have $m=n$.
        \begin{subproof}
            Fix $m,n$.
            Assume $m,n\in\dom{v}$ and $v(m)=v(n)$.
            Hence $m,n\in\naturals$.
            \begin{byCase}
                \caseOf{$m=n$.}
                    Follows.
                \caseOf{$m\neq n$.}
                    $m$ and $n$ are ordinals by
                        \obv{natural_number_is_ordinal}.
                    $m\in n$ or $n\in m$ by
                        \obv{ordinal_elem_connex}.
                    \begin{byCase}
                        \caseOf{$m\in n$.}
                            $n\in D$ by \obv{subseteq}.
                            $n\in C$ by \obv{subseteq}.
                            Hence for every $j\in n$ we have
                                $R(n)\subseteq
                                    \graphNeighbourhood{E}{v(j)}{X}$.
                            Hence $R(n)\subseteq
                                \graphNeighbourhood{E}{v(m)}{X}$.
                            $v(n)\in\edmGoodVertices{E}{R(n)}{\kappa}$.
                            Hence $v(n)\in R(n)$ by
                                \auto{edm_good_vertices_iff}.
                            Hence $v(n)\in
                                \graphNeighbourhood{E}{v(m)}{X}$ by
                                \obv{subseteq}.
                            $E$ is irreflexive by
                                \obv{simple_graph_irreflexive}.
                            Hence $v(m)\neq v(n)$ by
                                \obv{graph_neighbourhood_distinct}.
                            Contradiction.
                        \caseOf{$m\notin n$.}
                            Hence $n\in m$.
                            $m\in D$ by \obv{subseteq}.
                            $m\in C$ by \obv{subseteq}.
                            Hence for every $j\in m$ we have
                                $R(m)\subseteq
                                    \graphNeighbourhood{E}{v(j)}{X}$.
                            Hence $R(m)\subseteq
                                \graphNeighbourhood{E}{v(n)}{X}$.
                            $v(m)\in\edmGoodVertices{E}{R(m)}{\kappa}$.
                            Hence $v(m)\in R(m)$ by
                                \auto{edm_good_vertices_iff}.
                            Hence $v(m)\in
                                \graphNeighbourhood{E}{v(n)}{X}$ by
                                \obv{subseteq}.
                            $E$ is irreflexive by
                                \obv{simple_graph_irreflexive}.
                            Hence $v(n)\neq v(m)$ by
                                \obv{graph_neighbourhood_distinct}.
                            Contradiction.
                    \end{byCase}
            \end{byCase}
        \end{subproof}
        Follows by \auto{injective_function}.
    \end{subproof}
    Hence $v\in\bijections{\dom{v}}{\ran{v}}$ by
        \obv{injective_function_bijection_to_range}.
    Hence $v\in\bijections{\naturals}{\ran{v}}$.

    Let $K=\ran{v}$.
    $v\in\bijections{\naturals}{K}$.
    Hence $K$ is countably infinite by
        \auto{countably_infinite_iff_natural_bijection}.
    $K\subseteq X$ by \auto{funs_ran}.
    Show that for all $x,y\in K$ such that $x\neq y$ we have
        $x\mathrel{E}y$.
    \begin{subproof}
        Fix $x,y\in K$.
        Assume $x\neq y$.
        Take $m\in\dom{v}$ such that $v(m)=x$ by
            \obv{fun_ran_iff}.
        Take $n\in\dom{v}$ such that $v(n)=y$ by
            \obv{fun_ran_iff}.
        Hence $m,n\in\naturals$ and $m\neq n$.
        $m$ and $n$ are ordinals by \obv{natural_number_is_ordinal}.
        $m\in n$ or $n\in m$ by \obv{ordinal_elem_connex}.
        \begin{byCase}
            \caseOf{$m\in n$.}
                $n\in D$ by \obv{subseteq}.
                $n\in C$ by \obv{subseteq}.
                Hence for every $j\in n$ we have
                    $R(n)\subseteq\graphNeighbourhood{E}{v(j)}{X}$.
                Hence $R(n)\subseteq\graphNeighbourhood{E}{v(m)}{X}$.
                $v(n)\in\edmGoodVertices{E}{R(n)}{\kappa}$.
                Hence $v(n)\in R(n)$ by
                    \auto{edm_good_vertices_iff}.
                Hence $v(n)\in\graphNeighbourhood{E}{v(m)}{X}$ by
                    \obv{subseteq}.
                Hence $v(m)\mathrel{E}v(n)$ by
                    \obv{graph_neighbourhood_edge}.
            \caseOf{$m\notin n$.}
                Hence $n\in m$ by \obv{ordinal_elem_connex}.
                $m\in D$ by \obv{subseteq}.
                $m\in C$ by \obv{subseteq}.
                Hence for every $j\in m$ we have
                    $R(m)\subseteq\graphNeighbourhood{E}{v(j)}{X}$.
                Hence $R(m)\subseteq\graphNeighbourhood{E}{v(n)}{X}$.
                $v(m)\in\edmGoodVertices{E}{R(m)}{\kappa}$.
                Hence $v(m)\in R(m)$ by
                    \auto{edm_good_vertices_iff}.
                Hence $v(m)\in\graphNeighbourhood{E}{v(n)}{X}$ by
                    \obv{subseteq}.
                Hence $v(n)\mathrel{E}v(m)$ by
                    \obv{graph_neighbourhood_edge}.
                Hence $v(m)\mathrel{E}v(n)$ by
                    \obv{simple_graph_edge_symmetric}.
        \end{byCase}
    \end{subproof}
    Hence $K$ is a clique in $X$ with respect to $E$ by
        \obv{clique_in_graph}.
\end{proof}
